\documentclass[10pt,a4paper]{amsart}
\usepackage[margin=19mm]{geometry}
\usepackage{amsmath,amssymb,mathtools}
\usepackage[scr=rsfs,cal=euler]{mathalfa}
\usepackage{xfrac}
\usepackage[unicode,hidelinks,bookmarks=false]{hyperref}
\newcommand{\ie}{i.e.\ }
\newcommand{\cf}{cf.\ }
\newcommand{\resp}{respectively\ }
\newcommand{\Subsection}{\subsection}
\providecommand{\nobreakdash}{\nobreak\hbox{-}}
\setlength{\emergencystretch}{2em}
\multlinegap0pt
\usepackage{algorithm}\usepackage{algpseudocode}
\usepackage{amssymb}
\usepackage[shortlabels]{enumitem}
\usepackage[normalem]{ulem}
\newtheorem{lemma}{Lemma}
\renewcommand{\tt} {\mathtt}
\title{Universal cycle constructions for $k$-subsets and $k$-multisets}
\author{Colin Campbell, Luke Janik-Jones, Joe Sawada}
\date{Source: arXiv:2603.11954v1 (CC BY 4.0)}
\begin{document}
\maketitle

\noindent\emph{Source and licence.} Universal cycle constructions for $k$-subsets and $k$-multisets. Version arXiv:2603.11954v1 (CC BY 4.0), distributed by arXiv under the Creative Commons Attribution 4.0 International licence, http://creativecommons.org/licenses/by/4.0/, as recorded in the arXiv licence metadata for this exact version and shown on the abstract page https://arxiv.org/abs/2603.11954v1. Full licence notice: \texttt{licenses/arxiv-2603.11954-CC-BY-4.0.txt}.\par\smallskip

\noindent\textit{Editorial scope.} Counted unit: the article's lemma lem:props (source0.tex lines 734--743). The article's complete original proof of it is delivered with it (source0.tex lines 745--752, one \texttt{proof} environment of 8 lines). No mathematics is written, repaired or re-derived: the statement and the proof are the article's own lines, in the article's own notation, and no retained line is substituted or re-flowed.  No further source-local context is needed: every cross-reference of every retained line resolves inside the two delivered spans.  Nothing was omitted from any retained span, so the package carries no omission record.  No retained line contains a citation, so the wrapper carries no bibliography and no bibliography entry.  No retained line contains a figure environment, an image-inclusion command or drawing code, so no image is delivered and the package declares no asset dependency.  Editorial formatting only, in addition: an AMS \texttt{amsart} preamble that restates the article's own declarations and macros used by the retained lines, namely the environments \texttt{lemma} on separate counters, together with the standard packages those lines need.  The retained environments print in this document as Lemma 1 in order of appearance, whereas the article prints the same items with the numbering of its own full document.  The complete native source of the article as distributed in the arXiv e-print for this exact version is bundled unchanged as \texttt{sources/196246/source0.tex}.
\begin{lemma} \label{lem:props}
Let $\mathbb{T}$ be a PCR-based cycle-joining tree induced by the  first non-zero parent rule and let $\alpha = \tt{0}^{i-1}\tt{x}\tt{a}_{i+1}\cdots \tt{a}_n$ and
$\beta = \tt{0}^{j-1}\tt{y}\tt{b}_{j+1}\cdots \tt{b}_n$ be two distinct nodes (necklaces) in $\mathbb{T}$ where $i,j \leq n$ and $\tt{x},\tt{y} > 0$.  Then the following hold. 
\begin{enumerate}

    \item If $\alpha \neq \tt{0}^{n-1}\tt{1}$ then the parent of $\alpha$ is an aperiodic necklace. (All non-root periodic nodes are leaves.)
        \item $\mathbb{T}$ has the Chain Property.
    \item If $\alpha$ and $\beta$ are periodic then they have different parents.
\end{enumerate}
\end{lemma}

\begin{proof}
{\bf (1)} Since $\alpha$ is a necklace, its prefix $\tt{a}_1\cdots \tt{a}_i = \tt{0}^{i-1}\tt{a}_i$ is lexicographically less than or equal to every other substring in $\alpha$ of the same length when $\alpha$ is considered cyclically. Its parent $\gamma$ has prefix $\tt{0}^{i-1}(\tt{a}_i-1)$
which is strictly less than all other substrings in $\gamma$ considered cyclically, which means $\gamma$ is a necklace. Furthermore, this also implies $\gamma$ is aperiodic as long as $\gamma \neq \tt{0}^n$, which is the case only when $\alpha = \tt{0}^{n-1}1$.

{\bf (2-3)}
Suppose $\alpha$ and $\beta$ have the same parent $\gamma$.  If $i=j$ then $\alpha = \beta$, a contradiction to $\alpha$ and $\beta$ being distinct. Since they are necklaces where $\tt{x},\tt{y} > 0$, we have $\tt{a}_n, \tt{b}_n > \tt{0}$. Without loss of generality, suppose $i < j$.  The strings in the conjugate pair joining $\alpha$ and $\gamma$ have suffix $\tt{a}_n\tt{0}^{i-1}$ and the strings in the conjugate pair joining $\beta$ and $\gamma$ have suffix $\tt{b}_n\tt{0}^{j-1}$.  Thus, the Chain Property is satisfied. Furthermore, suppose $\alpha$ and $\beta$ are both periodic. Since $\beta$ is periodic, it must be that $\gamma$ has a substring $0^{j-1}$ appearing after the $j$-th symbol.  But this implies that $\alpha$ also has the same substring since  $\alpha$ and $\beta$ share the same parent. This contradicts the fact that $\alpha$ is a necklace.
Thus $\alpha$ and $\beta$ have different parents.  
\end{proof}

\end{document}
